% Research module: real-arithmetic bounds. Floating-point validation is separate.
\subsection{An enclosure through training-gate changes}
This section states the real-arithmetic stability inequality used by the
interval validator. A successful floating-point run is not by
itself a certificate; its input and rounding enclosures must also be verified.

Write $F(\theta)=J_\theta^*(x-f_\theta)$ for the selected vector field.
Fix a reference state $y$ and a radius $R$. Assume that along every straight
parameter segment within this ball, the smooth pieces obey
\[
 \max_i\|\E_n[e_\theta(x)x^T\alpha_i(x)]\|\le a.
\]
For each neuron and sample choose an upper bound
$h_{in}\ge N^{-1}\|x_n\|(w_i^Te_\theta(x_n))_+$ throughout the ball.
Set $d_{in}=|u_i(y)^Tx_n|/\|x_n\|$ for nonzero samples, and
\[
 B_i(s)=\sum_{n:d_{in}\le s}h_{in},\qquad 0\le s\le R.
\]
If the center ranges along a reference subinterval, lower bounds for the
margins and uniform upper bounds for the weights can replace these quantities.

\begin{lemma}[One-sided stability with adverse gate jumps]\label{initcert:tube}
Suppose $B_i(s)\le a_i+\beta s$ for $0\le s\le R$, where $a_i,\beta\ge0$.
Then for $\|z\|\le R$,
\[
 \langle z,F(y+z)-F(y)\rangle
 \le(a+\beta)\|z\|^2+\|(a_i)_i\|_2\|z\|.
\]
Compatible endpoint selectors are allowed. Along an absolutely continuous
reference path $y(t)$ with defect $\|F(y(t))-\dot y(t)\|\le\eta$, the
error $r(t)=\|\theta(t)-y(t)\|$ therefore satisfies, while $r\le R$,
\[
 D^+r\le(a+\beta)r+\|(a_i)_i\|_2+\eta
\]
in the comparison sense (or almost everywhere away from $r=0$).
\end{lemma}
\begin{proof}
On a smooth activation region,
$DF=-J^*J+H_e$. The residual Hessian $H_e$ is block diagonal across
neurons, with each block having off-diagonal entries $A_i,A_i^T$, where
$A_i=\E_n[e x^T\alpha_i]$. Thus $\lambda_{\max}(DF)\le a$.
Integrate this inequality along the segment $y+sz$, separating its finitely
many training hyperplane crossings. When one gate crosses, the vector-field
jump has only an input-coordinate component
$\Delta F_{u_i}=N^{-1}\Delta\alpha_i(w_i^Te_x)x$; the output-coordinate
component is continuous because the preactivation is zero there. Its pairing
with $z$ is at most
$N^{-1}(w_i^Te_x)_+|z_{u_i}^Tx|\le h_{in}\|z_{u_i}\|$.
Only samples with $d_{in}\le\|z_{u_i}\|$ can cross. The sum of jumps is
therefore bounded by
\[
 \sum_i\|z_{u_i}\|B_i(\|z_{u_i}\|)
 \le\|(a_i)_i\|_2\|z\|+\beta\|z\|^2.
\]
Simultaneous crossings add; endpoint partial selectors satisfy the same
upper bound by taking partial jumps. Segments that lie in a hyperplane
contribute no conormal pairing for that gate. This proves the first claim.
Take the inner product of the error equation with $\theta-y$ and divide
by $r$ when it is positive. Standard norm regularization gives the comparison
at zero.
\end{proof}

For a finite set of distances, permissible intercepts are explicit:
\[
 a_i(\beta)=\max\big\{0,\max_{0\le s\le R}(B_i(s)-\beta s)\big\}.
\]
The inner maximum only needs the sorted sample distances, including zero.
There is no density assumption or asymptotic substitution for these counts.

\begin{corollary}[A verified step criterion]\label{initcert:tube-step}
On a reference interval of length $h$, let the preceding bounds be uniform
on its radius-$R$ tube. Put $k=a+\beta$ and
$\omega=\|(a_i)_i\|_2+\eta$. If the initial error is at most $r_0<R$ and
\[
 e^{kh}r_0+\omega\frac{e^{kh}-1}{k}<R,
\]
then the flow remains in the tube on the whole interval and its terminal
error is bounded by the left side. At $k=0$ the quotient is interpreted as
$h$. Successive verified steps give an enclosure from the initial state.
\end{corollary}
\begin{proof}
Apply the scalar integrating-factor inequality until the first exit. Its
upper bound is nondecreasing in time and strictly smaller than $R$ at the
terminal time, excluding a first exit. Iterate using the terminal upper bound
as the next initial radius.
\end{proof}
The absolute reference defect must also include any gates crossed by the
reference path. The favorable sign used in Lemma~\ref{initcert:tube} applies to
\emph{stability}, not automatically to the defect of a numerical interpolant.
Using only a smooth-region defect bound would omit this term and would not
constitute a certificate.

\subsubsection{Cubic-reference defect used in the local validator}
Let $y(s)$ be the cubic Hermite interpolant of two stored parameter states
and two stored velocities, over a step of exact length $1/5000$.
The stored velocities are merely coefficients of the reference polynomial;
they need not be exact field values. Put $a=1/10000$ and center the step
at $s=0$. Fix the actual training masks at $y(0)$ and denote that polynomial
field by $F^0$. The validator encloses
\begin{align*}
 d_0&=\|F^0(y)-y'\|,\\
 d_1&=\|DF^0(y)y'-y''\|,\\
 d_2&=\|D^2F^0(y)[y',y']+DF^0(y)y''-y'''\|
\end{align*}
at the midpoint. Suppose $P_*,V_*,A_*,T_*$ bound respectively
$\|y\|,\|y'\|,\|y''\|,\|y'''\|$ throughout the step, and $L_*$
bounds $\|DF^0\|$. For the covariance bound $\lambda$,
\[
 \|D^2F^0\|\le3\lambda P_*,\qquad
 \|D^3F^0\|\le3\lambda.
\]
Indeed the fixed-mask output is quadratic with second derivative norm
at most $\sqrt\lambda$, while $\|J\|\le\sqrt\lambda P_*$;
differentiate $F^0=J^*e$ twice and three times.
Taylor's integral remainder therefore gives
\[
 \eta_{\rm sm}=d_0+ad_1+\tfrac12a^2d_2+
 \tfrac16a^3\big(3\lambda V_*^3+9\lambda P_*V_*A_*+L_*T_*\big)
\]
as an upper bound for the fixed-mask defect on the whole step.
The third derivative of the defect uses $y''''=0$.

The true defect additionally includes $\eta_{\rm gate}$, the difference
between the actual field and $F^0$ along the reference polynomial.
It is bounded by the direct field argument in
Proposition~\ref{initcert:static-bound}, using neuronwise input and output excursions
\[
 d_i^u=a\|y'_{u_i}(0)\|+\tfrac12a^2\|y''_{u_i}(0)\|
                    +\tfrac16a^3\|y'''_{u_i}\|,
\]
and the analogous $d_i^w$. The corresponding whole-parameter excursion
bounds residual changes. Thus $\eta=\eta_{\rm sm}+\eta_{\rm gate}$ is
uniform on the step, including all reference-path gate changes.

For the one-sided rate on a tube, let $A_i^0$ be its midpoint residual
matrices and $D_*$ the pointwise output-change coefficient from the midpoint
throughout the tube. If $p_{in}$ flags every potentially changed gate there,
a valid smooth-piece bound is
\[
 a_* = \max_i\|A_i^0\|+\lambda D_*
       +\max_i N^{-1}\sum_n p_{in}\|e_y(x_n)\|\|x_n\|.
\]
For gate jumps the validator bounds the positive conormal residual throughout
the same tube, sorts lower bounds on its reference-path gate distances,
and sums upward-rounded weights. With $\beta=1$, the intercept
$\max_s(B_i(s)-s)_+$ only requires those finite sorted distances.
All tube radii are explicit inputs, rather than inferred from successful
numerical integration. Corollary~\ref{initcert:tube-step} is checked at every
step. Its strict inequality is what validates the reference, independently
of the integration method originally used to choose its coefficients.
